% !TeX program = xelatex
% CurveLab - standalone vector export. Compile with XeLaTeX, not pdfLaTeX.
% Structure: tailieuvehinh.pdf, I.1 (p.4); coordinates, draw, node, scope, clip.
% 3D is the current parallel projection, with visible/hidden line parts (I.3.5, pp.30-32).
% Captures the current viewport and sampled paths, without embedding a PNG.
% Width 16 cm; line widths and label sizes preserve the current drawing proportions.
% Mode: Primitive test
\documentclass[varwidth=50cm,border=2pt]{standalone}
\usepackage{fontspec}
\IfFontExistsTF{Times New Roman}{\setmainfont{Times New Roman}}{\setmainfont{TeX Gyre Termes}}
\IfFontExistsTF{Arial}{\setsansfont{Arial}}{\setsansfont{TeX Gyre Heros}}
\usepackage{amsmath,amssymb}
\usepackage{tikz,graphicx}
\usetikzlibrary{calc,angles,arrows.meta}
\definecolor{cl0}{RGB}{255,255,255}
\definecolor{cl1}{RGB}{20,99,214}
\definecolor{cl2}{RGB}{255,169,65}
\definecolor{cl3}{RGB}{0,0,0}
% Change only this width; geometry, strokes and labels scale together.
\newcommand{\FigureWidth}{16cm}
\begin{document}
\begin{center}
\resizebox{\FigureWidth}{!}{%
\begin{tikzpicture}[x=1cm,y=-1cm,
 declare function={figwidth=16; figheight=10;},
  s1/.style={fill=cl0,fill opacity=1},
  s2/.style={draw=cl1,draw opacity=1,line width=1.13811pt,line cap=butt,line join=miter},
  s3/.style={fill=cl2,fill opacity=0.45},
  s4/.style={draw=cl1,draw opacity=1,line width=1.13811pt,line cap=butt,line join=miter,dash pattern=on 2.845276pt off 2.27622pt},
  s5/.style={anchor=base west,inner sep=0pt,outer sep=0pt,text=cl3,text opacity=1,font={\rmfamily\bfseries\fontsize{10.242992}{12.291591}\selectfont}},
  s6/.style={anchor=north east,inner sep=3pt,fill=white,fill opacity=.94,text opacity=1,font={\fontsize{8.535827}{10.242992}\selectfont}}
]
% Coordinates use the displayed projection; original coordinates are retained in comments.
\path[use as bounding box] (0,0) rectangle ({figwidth},{figheight});
\begin{scope}
\clip (0,0) rectangle ({figwidth},{figheight});

% Drawing in display order: axes/grid, shapes, labels, formula.
\fill[s1] (0,0)--(16,0)--(16,10)--(0,10)--cycle;
\draw[s2] (0,4)--(16,4);
\begin{scope}
\clip (2,2) rectangle (12,7);
\fill[s3] (5.4,2)--(12,2.694737)--(12,3.84)--(8.961538,7)--(7,7)--(2,2)--(2,2)--cycle;
\end{scope}
\draw[s4] (7,5.2) ellipse [x radius=1.8cm,y radius=1.8cm];
\node[s5] at (4,1.2) {Điểm A₁: 90° \& 50\% \_ \{x\}};
% Formula shown in the selected corner.
\node[s6] at (15.72,0.28) {$\begin{gathered}x(t)=\cos t\\y(t)=\sin t\end{gathered}$};
\end{scope}
\end{tikzpicture}%
}
\end{center}
\end{document}
